\section{A single product equation with 111 arithmetic operations}
\label{sec:product-equation}
The three-root norm can be replaced by one equality between products.
The price is a higher total degree. The resulting two equations have
eleven unrestricted integer variables, using the same entire source-34
coefficient tuple as Section~\ref{sec:short-carrier}:
\[
\begin{array}{l|rr}
\text{first-index modulus}&\text{degree}&\text{operations}\\\hline
XY&966&111\\
4q^4&882&112.
\end{array}
\]
Both improve the inspected all-root fusion equation of degree 1632 in
114 operations. The degree-440 construction remains the lower-degree
alternative at a greater arithmetic cost.

\subsection{Definition of the two polynomials}
Keep the eleven inputs $(j,C,z,h,r,v,w,T,c,n,s)\in\Z^{11}$ and the
common definitions through $P=D_\tau D_g$ in
Section~\ref{sec:short-carrier}. For the 111-operation equation use
$k=R+1+hXY$; for the 112-operation equation use $k=R+1+4hq^4$.
Thus, in both cases,
\[
 D_\tau=1+XY^2(XY^2+1)k^2,
 \qquad D_g=1+4(3\mathcal B-4)D_\tau s^2.
\]
Replace the strong and auxiliary definitions by
\begin{align*}
 A&=a+2,&D_m&=1+\Delta c^2,&U&=rc^2 A D_m,\\
 D_f&=1+\Delta U^2P,&Q&=\Delta^2U^2P,&\delta&=Q-1,\\
 b_0&=c-R+4UPT,&V&=D_fb_0-c,&C_y&=QV^2-1.
\end{align*}
Use the same short projector and scalar carrier,
\begin{align*}
 K&=X(3c-2X)+2,&S&=srH,&z_c&=HT+K,\\
 \mathcal A&=(z_c^2+3H^2)(z_c^3-19H^3),&&
 Z=Sn-\mathcal A.
\end{align*}
The requested equation is simply
\begin{equation}\label{eq:single-product}
 \boxed{\delta Z^2-S^2P D_f C_y=0.}
\end{equation}
All abbreviations are polynomials and introduce no extra variables.
In the literal evaluator,
\[
 P=D_\tau+(3\mathcal B-4)(2D_\tau s)^2,
\]
The short-index graph reuses $4q^2$ from the budget to form
$(4q^2)q^2=4q^4$, saving one more multiplication. The product $UP$ is shared between
$\Delta U^2P=(\Delta U)(UP)$ and $b_0=c-R+4(UP)T$.
The complete graphs are \code{lowops/product11\_degree966.json} and
\code{lowops/product11\_degree882.json}. Each is bound to all six original
fixed numerals; no coefficient parameter remains free.

\begin{theorem}\label{thm:single-product}
Each polynomial in~\eqref{eq:single-product} has an integer zero exactly
when the retained native source has a finite refutation. Composing its
ordinary source calibration gives
\[
 \ZFC\vdash
 \bigl(\text{no zero of the specified polynomial in }\Z^{11}\bigr)
 \longleftrightarrow\Con(\ZFC).
\]
The exact degrees and arithmetic counts are the two rows displayed above.
\end{theorem}

\subsection{Coprimality recovers the signs as well as the squares}
The unconditional affine bounds from the preceding section give
$A\ne0$, $H\ne0$ odd, $\Delta=A^2-1>1$, and $X\equiv0\pmod8$.
For every integer $\ell$, $\ell(\ell+1)\ge0$. Therefore
\[
 D_\tau\ge1,\quad D_\tau\equiv1\pmod8,\quad
 D_g\equiv1\pmod4,\quad P\equiv1\pmod4.
\]
In particular $P\ne0$. Neither $D_f$ nor $\delta$ can vanish, because
$\Delta>1$. Also $C_y\ne0$: the integer identity $QV^2=1$ would force
$Q=1$, which is impossible. Put $D_a=\delta C_y$; it too is nonzero.
These statements do not assume that $P,Q$, or $D_f$ are positive.

If $S=0$, the polynomial becomes $\delta\mathcal A^2$, which is nonzero.
Indeed
\[
 \mathcal A=H^5\pi(T+K/H),\qquad
 \pi(t)=(t^2+3)(t^3-19),
\]
and $\pi$ has no rational root. Thus a zero has $S\ne0$, so $s,r\ne0$.
We deliberately leave $c=0$ open until the guard is recovered.

The three integers $P,D_f,D_a$ are pairwise coprime. The first two
congruence computations are immediate:
\[
 D_f\equiv1\pmod P,\qquad D_a\equiv1\pmod P.
\]
For the third, reduction modulo $D_f$ gives
\[
 Q\equiv-\Delta,\quad\delta\equiv-A^2,\quad
 V\equiv-c,\quad C_y\equiv-D_m,\quad D_a\equiv A^2D_m.
\]
But $AD_m\mid U$, so $D_f\equiv1$ modulo both $A$ and $D_m$.
Consequently $\gcd(D_f,D_a)=1$. All gcds refer to absolute values;
the computations hold at negative inputs and at $U=0$.

Multiply~\eqref{eq:single-product} by $\delta$ and divide by $S^2$:
\begin{equation}\label{eq:coprime-square-product}
 \left(\frac{\delta Z}{S}\right)^2=P D_fD_a.
\end{equation}
A rational square root of an integer is integral. Pairwise coprimality
therefore makes $|P|,|D_f|,|D_a|$ separately squares. Since
$P\equiv1\pmod4$, it cannot be a negative odd square. Hence $P>0$,
then $D_f=1+\Delta U^2P>0$, and then~\eqref{eq:coprime-square-product}
forces $D_a>0$. There are positive integers
\[
 p_0^2=P,\qquad f^2=D_f,\qquad e_0^2=D_a.
\]
The further identity $\gcd(D_\tau,D_g)=1$ splits $P$ once more:
\[
 D_\tau=\tau^2,\qquad D_g=g^2,\qquad p_0=\tau g.
\]
Because $s\ne0$, the inequality $\mathcal B\le1$ would imply
$D_g\le1-4D_\tau<0$. Thus $\mathcal B\ge2$. The centered budget now
gives the common source signs and $|c/k-Y|<1/2$. It also excludes $c=0$:
at that value,
\[
 \mathcal B=k^2\bigl(q^2-D_{\rm err}-4q^2Y^2\bigr)\le0.
\]
We have obtained $c\ne0$ and $U\ne0$ without an initial assumption about
the signs of the saturated radicands.

\subsection{Recovering the original Pell witnesses}
Define, only within the proof,
\[
 I=Up_0,\qquad T_{\rm old}=p_0T.
\]
These are integers, $I\ne0$, $c^2\mid I$, and $|I|\ge c^2$.
The restored equalities are exact:
\begin{align*}
 f^2&=1+\Delta I^2,&Q&=(\Delta I)^2>1,\\
 4IT_{\rm old}&=4UPT,&
 V&=f^2(c-R+4IT_{\rm old})-c.
\end{align*}
Moreover $y=e_0/\delta\in\Q$ satisfies
$QV^2-(Q-1)y^2=1$. The adjacent-coefficient descent used in
Section~\ref{sec:short-carrier} makes $y$ integral. Its hypotheses
$Q>1$ and integer $V$ have now been established. Cancellation of the
nonzero $\delta$ in~\eqref{eq:single-product} yields
\[
 Z^2=S^2p_0^2f^2y^2,
\]
so $Z/S$ is an integer. Hence $S\mid\mathcal A$ and $H\mid K^5$.

The complete signed folding and rank argument now has its original
hypotheses: the budget strip, the first Pell square, a nonzero strong
ordinate divisible by $c^2$, and the auxiliary norm with
$V\equiv-c\pmod f$, $V\equiv-R\pmod c$. Its use does not require
an odd strong multiplier. In the $4q^4$ version the first-index argument
is precisely the one already proved in Section~\ref{sec:short-carrier}.
The exact index and orientation later restore
$XY\mid k-R-1$, so the original quotient $h$ can be regenerated.

The fifth-power height argument also applies unchanged. With
$R=2t+1$, $t\ge24$, $x=|X|\ge2^{64}$ and
\[
 \Pi=(\operatorname{sign}(c)X-2^R)
          (2^R\operatorname{sign}(c)X+1),
\]
it gives $H\mid\Pi^5$ and $|\Pi|^5<|H|$. Thus $\Pi=0$, the exact
main projection is restored, and the signed half-binomial argument
recovers every outer native condition. Fresh private witnesses supply
an actual zero of the retained source. This proves the forward
finite-refutation implication using the same source calibration.

\subsection{Completeness and the extra interpolation factor}
Start with a canonical outer source witness. It has $A>1$ even,
$R\equiv3\pmod4$, $c=\psi_A(R)>0$, $c\equiv R\pmod4$,
$D_\tau=\tau^2$, and the stronger budget $\mathcal B\ge64$.
The positive nonsquare $(3\mathcal B-4)D_\tau\equiv2\pmod3$
has a Pell unit. Squaring one supplies positive $s,g$ with
$g^2=1+4(3\mathcal B-4)D_\tau s^2$. Fix them and put $p_0=\tau g$.

Choose one positive strong pair $I,f$ satisfying
\[
 4sAD_m p_0c^2\mid I,
 \qquad f^2-\Delta I^2=1.
\]
Set $r=I/(AD_mp_0c^2)$ and $U=I/p_0$. The strong pair is now fixed.
Every prime factor of $M=p_0sr$ divides $I$. The canonical projector
gives $\kappa=K/H\in\Z$. Choose an intersective residue $a_0$ modulo
$sr$ with $\pi(a_0+\kappa)\equiv0\pmod{sr}$. Supported-modulus
interpolation, at this same fixed $I,f$, supplies the full auxiliary
witness with
\[
 T_{\rm old}\equiv p_0a_0\pmod{p_0sr}.
\]
The extra factor $p_0$ in this modulus is essential. It makes
$T=T_{\rm old}/p_0$ an integer and ensures $T\equiv a_0\pmod{sr}$.
Therefore
\[
 \frac{\mathcal A}{S}
     =\frac{H^4\pi(T+\kappa)}{sr}\in\Z.
\]
Finally set $n=\mathcal A/S+p_0fy$. Equation~\eqref{eq:single-product}
vanishes. The order of choices is outer witness, $s,g$, one strong
$I,f,r$, a residue, $T_{\rm old},T,y$, and $n$. No coprimality between
these moduli is assumed, and the strong pair is never changed after
the residue request. This proves the reverse implication.

\subsection{Exact degree and finite validation}
The key degrees in the two literal graphs are
\[
\begin{array}{c|rrrrrrrrrr}
 &k&P&U&D_f&Q&b_0&V&C_y&S&\mathcal A\\\hline
XY&7&94&23&152&164&118&270&704&8&35\\
4q^4&5&82&23&140&152&106&246&644&8&35.
\end{array}
\]
In both cases $\deg Z=35$. The summand $\delta Z^2$ has degree 234
or 222, while $S^2PD_fC_y$ has degree 966 or 882. The latter therefore
determines the full leading homogeneous part.

Put $q_0=6\beta j$, $X_0=12q_0w$, $Y_0=2vq_0^3$,
$a_0=X_0Y_0$, and
\[
 k_0=ha_0\quad\hbox{or}\quad
 k_0=4hq_0^4-q_0^3(12wC-q_0z),
\]
respectively. Direct exact multiplication gives the leading form
\[
 2^{36}3^7q_0^{14}X_0^{70}Y_0^{112}k_0^{42}
       s^{16}r^{12}c^{40}T^2.
\]
For the native modulus its coefficient at
$j^{588}w^{112}v^{154}h^{42}s^{16}r^{12}c^{40}T^2$ is
$2^{1002}3^{707}\beta^{588}>0$. For the short modulus the coefficient
at $j^{588}w^{70}v^{112}h^{42}s^{16}r^{12}c^{40}T^2$ is
$2^{960}3^{665}\beta^{588}>0$. Their exponent sums give the two
stated exact degrees. The latter coefficient is selected by the maximal
$h$ power, so no cancellation from the other leading summands is possible.

The independent checker uses only the Python standard library and
does not import the producer. It verifies every local polynomial formula,
all 135 or 136 live nodes, all six complete coefficient payloads,
the coprimality remainder identities, both saturation identities,
the new index-modulus identity, and exact leading monomials at every
gate. Five altered mathematical formulas in each graph are rejected.
Normal and optimized Python runs agree, and both graphs rebuild
byte-for-byte. These finite checks certify the exported algebra and
counts; the unbounded semantic argument is the proof above and its
explicitly retained source lemmas.
